% This is samplepaper.tex, a sample chapter demonstrating the
% LLNCS macro package for Springer Computer Science proceedings;
% Version 2.21 of 2022/01/12
%
\documentclass[runningheads]{llncs}
%
\usepackage{subcaption}
\usepackage[T1]{fontenc}
% T1 fonts will be used to generate the final print and online PDFs,
% so please use T1 fonts in your manuscript whenever possible.
% Other font encondings may result in incorrect characters.
%
\usepackage{graphicx}
\usepackage{amsmath}
\usepackage{booktabs}
\usepackage{float}
\usepackage{multirow}
\usepackage{amsfonts}
\usepackage{xcolor}
\usepackage{amssymb}
% Used for displaying a sample figure. If possible, figure files should
% be included in EPS format.
%
% If you use the hyperref package, please uncomment the following two lines
% to display URLs in blue roman font according to Springer's eBook style:
%\usepackage{color}
%\renewcommand\UrlFont{\color{blue}\rmfamily}
%


%\title{Learning and Recognizing Latent Innovation Maturity Indicator Patterns in Weakly Supervised Text}
%GL: et pourquoi pas simplement :
\title{Learning and Recognizing Latent Innovation Maturity Indicator Patterns in Texts}
%GL: dans la recherche du titre optimal peut être indiquer un truc sur la méthode en intégrant quelque chose comme : Weakly Supervised Technology Readiness Inference from textual data 

%
\titlerunning{Latent TRL Pattern Recognition}
% If the paper title is too long for the running head, you can set
% an abbreviated paper title here
%

\author{Mélusine Caillard\inst{1,2,3} \and
Gaël Lejeune\inst{1}\orcidID{0000-0002-4795-2362} \and \\ 
Pierre-Emmanuel Fayemi\inst{3} \and
Aoussat Améziane\inst{2}}
%
\authorrunning{M. Caillard et al.}
% First names are abbreviated in the running head.
% If there are more than two authors, 'et al.' is used.
%
\institute{STIH, Sorbonne Université, Paris \and
LCPI, ENSAM, Paris \and IKOS Lab, Levallois-Perret}
%
\begin{document}
\maketitle 

% typeset the header of the contribution
%
\begin{abstract}
Technology Readiness Levels (TRLs) are an ordinal  maturity scale widely used in innovation assessment, yet inferring them from text is hindered by scarce labeled data and domain-dependent conventions. We formulate TRL inference as a latent ordinal pattern recognition problem under weak supervision. We introduce a large silver corpus  constructed from project descriptions using hybrid symbolic–semantic labeling, and reconstruction rules to improve alignment with expert annotations. 
We compare flat classification, regression, ordinal decomposition, and hierarchical models using ordinal error metrics. Linear models recover a coherent ordinal pattern in-domain, with errors concentrated within $\pm$1 level. We show that this structure degrades under zero-shot transfer to an expert gold corpus, revealing a strong semantic domain shift. We also show that rigid ordinal constraints are brittle under transfer, while soft hierarchical routing and lightweight hybrid ensembling significantly reduce large ordinal errors. Analyses of grammatical and structural representations indicate that domain shift is driven primarily by semantic conventions rather than syntax. These results highlight both the potential and the limits of weakly supervised ordinal pattern recognition in real-world settings. To foster reproducibility, we release the dataset and the code in a public GitHub repository.


\keywords{Text mining  \and Weak supervision \and Pattern recognition.}
\end{abstract}
%
%
%
\section{Introduction}

In an era of accelerated technological development, the ability to reliably assess
the maturity of emerging technologies is critical for public funding agencies,
private investors, and R\&D managers. To standardize this assessment, the
Technology Readiness Level (TRL) scale, originally introduced by NASA
~\cite{mankins1995}, later became an ISO standard and adopted by the European Union, has become a
widely used benchmark \cite{Heder2017TRL}. 
The TRL scale ranges from TRL~1 (basic principles observed) to TRL~9 (system
proven in operational environment). It provides a common language to map
innovation pathways, and to identify critical transitions such as the move from
laboratory validation (TRL~4) to industrial demonstration (TRL~6--7). However,
applying this metric consistently to massive datasets like portfolios of research projects remains a significant challenge.
TRLs are well-defined as an expert scale but detecting them
in unstructured project descriptions is non-trivial.
The maturity of a technology is rarely stated explicitly; rather, it emerges as a
\emph{latent semantic pattern} distributed across the choice of verbs
(e.g., ``investigate'' vs. ``deploy''), the level of technical detail, and the
description of the operating environment.

Current automated approaches are severely constrained by \emph{data scarcity}~\cite{Voltan2024} and, to our knowledge, there is no large annotated corpus for the TRL assessment task. Constructing such a resource would require domain experts to annotate thousands of technical summaries, which is both costly and time-consuming. As a result, most practical systems rely on manual assessment by expert teams. 
%\subsection{Contribution}
We propose a \emph{frugal and weakly
supervised} framework for TRL assessment in large textual corpora. Our
contributions are threefold. 
We design a hybrid labeling function that combines high-precision symbolic rules with high-recall semantic inference, automatically generating a large-scale ``Silver standard'' dataset ($N = 35{,}104$) from the CORDIS corpus 
%\footnote{The Community Research and Development Information Service, European Union’s database of research and innovation projects \url{https://cordis.europa.eu/datalab}}. %GL: on pourrait gratter sur la footnote qui fait 2 lignes :
\footnote{European Union’s database of research/innovation projects \url{cordis.europa.eu/datalab}}.
It provides comprehensive metadata together with public textual fields for all Horizon 
%2020 and Horizon Europe 
research projects. 
Because of its size, thematic diversity, and standardized structure, CORDIS constitutes a unique corpus for studying technological maturity at scale, despite the fact that explicit TRL mentions are rare.
We benchmark both frugal models and Transformer-based deep models (BERT base and SciBERT) on an independent expert-annotated “Gold” dataset derived from the IEA Clean Energy Technology Guide under zero-shot and few-shot transfer protocols. This public resource released by the International Energy Agency\footnote{\url{www.iea.org/data-and-statistics/data-tools/etp-clean-energy-technology-guide}}
 provides technology notices that include expert-validated TRL assessments, offering a high-quality benchmark for evaluating transfer to a real-world, domain-specific maturity taxonomy.
We analyse TRL patterns, studying (i) their ordinal structure, (ii) their sensitivity to domain shift, and (iii) the contribution of grammatical cues as a possible formalization of technological maturity.
The structure of the paper is as follows: Section 2 reviews related work, Section 3 presents the proposed framework, and Sections 4 and 5 report the experiments and results, respectively.

\section{Related Work on Technology Readiness Levels}

TRLs were originally introduced at NASA as a nine-point scale to characterise
the maturity of space technologies~\cite{mankins1995} and were later adopted in
European research policy. In practice, TRLs are most often assessed through
expert judgement, sometimes supported by structured questionnaires or
checklists~\cite{engel2012}. Several authors have proposed spreadsheet-based
``TRL calculators'' that guide experts through series of questions
\cite{nolte2008,nolte2003}, or calibrated TRL assignments against mission data
to link maturity to project cost milestones~\cite{terrile2015}. These tools
improve transparency and repeatability but remain fundamentally manual and
project-specific, and have been criticised for ambiguity and inconsistent
interpretation across organisations and sectors~\cite{olechowski2020}.

Beyond individual projects, matrix- or survey-based frameworks have been
developed to assess the readiness of technologies or portfolios at a higher
level. Chukhray et al.~\cite{chukhray2020} propose structured criteria and expert scoring for
assessing technology transfer readiness. Mortara and Ford~\cite{MortaraFord2012}
embed TRL-type assessments into broader decision tools for technology
acquisition and protection. These approaches support strategic decision-making
but require substantial expert input and do not scale to tens of
thousands of textual project descriptions.

%A complementary line of
Other works use publication, patent and citation data to
approximate technological maturity. Watts and Porter~\cite{wattsporter1997}
discuss innovation forecasting from bibliometric signals, while Campani and
Vaglio~\cite{campanivaglio2015} interpret the evolution of scientific and
technological impact through hype-type growth curves. Lezama-Nicolás et
al.~\cite{lezamanicolas2018} propose a bibliometric method based on publication
growth and citation patterns. Recent surveys such as \cite{loginova2025}
emphasise the potential of large text corpora for maturity assessment, but also
underline the lack of standardised, scalable methods. Furthermore, bibliometric
approaches typically operate at a meso-level (domains or technologies) rather
than assigning TRLs to documents.

Closer to our work, some studies have addressed document-level TRL prediction.
Britt et al.~\cite{britt2008} apply classification techniques to NASA
reports, showing that machine learning can recover maturity information on
relatively small, domain-specific corpora. Cunningham~\cite{cunningham2014}
discusses big data for TRL analysis but does not provide %a fully
an operational pipeline. More recently, Voltan et al.~\cite{Voltan2024} describe
%a methodology to build labelled corpora and classifiers for defence technologies, relying on expert annotation and supervised learning. 
a methodology to build labeled corpora with expert annotation to train classifiers designed for defence technologies. %, relying on expert annotation and supervised learning. 
These contributions demonstrate the feasibility of TRL classification from text but generally depend on manually labeled datasets, focus on a single sector, and
do not explicitly analyse domain shift.

Compared to prior work, our approach combines three ideas that %, to our knowledge, 
have not been studied jointly: (i) a weakly supervised pipeline to
construct a large-scale Silver-standard corpus of TRL-labeled project
descriptions; (ii) a systematic analysis of TRLs as a latent
ordinal pattern, including regression, classification and grammar-aware
representations; and (iii) an explicit evaluation of domain shift by
transferring from multi-sector CORDIS projects to an expert-annotated Gold
dataset (IEA Clean Energy Technology Guide). In this sense, our work
complements expert-, matrix- and bibliometric approaches by providing a
scalable, text-based framework that is both frugal (linear models on weak
labels) and adaptable (few-shot fine-tuning of specialised language models).

\section{Weak Supervision Framework}

Instead of relying on manual annotation, we construct a hybrid labeling function
that combines the high precision of symbolic rules with the high recall of
semantic inference. This pipeline automatically processes raw project
descriptions to generate a training corpus, hereafter referred to as the
\emph{Silver Dataset}.
To validate this framework, we adopt a realistic scenario: we train on abundant, noisy, administrative data (CORDIS), which covers multiple R\&D domains but largely \emph{excludes} the energy technologies %found in the IEA corpus, 
and we evaluate on scarce, high-quality, expert data focused on the energy domain (IEA) in order to evaluate the robustness of the approach. %GL: peut être ici un annonce du plan de la section
The remainder of this section details the hybrid labeling function and the resulting Silver corpus.
\subsection{Hybrid Labeling Function}
\label{sec:hybrid}
%MC : math formalization
\textcolor{purple}{We consider a text description $x \in \mathcal{X}$ and aim to predict a technological maturity level $y \in \{1,\dots,9\}$, which is naturally ordered. We assume the existence of a latent continuous score $z_\theta(x)\in\mathbb{R}$ together with ordered thresholds $\theta_0 < \theta_1 < \cdots < \theta_9$ such that
\(
y = k \quad \Longleftrightarrow \quad \theta_{k-1} < z_\theta(x) \le \theta_k.
\)
This formulation follows standard ordinal regression settings, where the observed TRL results from the discretization of a latent maturity signal expressed in text.}


\textcolor{blue}{The training data is not directly labeled with ground-truth TRL values, but obtained through a weak supervision process. We define a labeling function $\mathcal{F} : \mathcal{X} \rightarrow \{1,\dots,9\} \cup \{\varnothing\}$ such that }
\color{blue}
\[
\mathcal{F}(x) =
\begin{cases}
\lambda_{\text{sym}}(x) & \text{if } \lambda_{\text{sym}}(x) \neq \varnothing, \\
\lambda_{\text{sem}}(x) & \text{if } \max_j \cos(\phi(x),\phi(d_j)) \ge \tau, \\
\varnothing & \text{otherwise}.
\end{cases}
\]

This induces a weakly labeled dataset
\[
\tilde{\mathcal{D}} = \{(x_i, \tilde{y}_i)\}_{i=1}^N,
\]
where $\tilde{y}_i = \mathcal{F}(x_i)$ are noisy proxies of the latent true labels $y_i$.

The learning objective is therefore to estimate a predictor $z_\theta : \mathcal{X} \rightarrow \mathbb{R}$ by minimizing an ordinal loss under label noise:
\[
\min_\theta \ \mathbb{E}_{(x,\tilde{y}) \sim \tilde{\mathcal{D}}}
\big[\mathcal{L}_{\text{ord}}(z_\theta(x), \tilde{y})\big].
\]   

This setting relates to weak supervision and data programming approaches, in which heuristic labeling functions are used to generate large training sets without exhaustive manual annotation \cite{ratner2017snorkel}.

The resulting labels should therefore be viewed as structured noisy annotations rather than exact ground truth, a setting closely related to learning under label noise \cite{frenay2014labelnoise}.

\color{black}
%MC : end additional formalization

The symbolic heuristic, $\lambda_{\text{sym}}$, captures explicit TRL mentions. We
designed a set of regular expressions to identify patterns such as ``TRL~X'',
``Technology Readiness Level~X'', or ranges like ``TRL~m--n''. Because project
texts may contain several TRL mentions (for instance an initial and a target
TRL, and sometimes wide ranges), we adopt a \emph{conservative minimization
rule}: for any range $[y_{\text{start}}, y_{\text{end}}]$ we keep the lower
bound $y_{\text{start}}$, and if multiple TRL values are detected in the same
document we assign the minimum over all candidates. This yields a lower bound
on the claimed maturity and favours precision over optimistic reporting.
This component yields near-perfect precision but suffers from extremely low
recall, as explicit mentions are rare in administrative summaries.
For the semantic heuristic we use the TRL definitions published by the European Commission as semantic anchors
$D = \{d_1, \dots, d_9\}$. Both the input text $x$ and the definitions $d_k$
are mapped into a word $n$-grams TF--IDF vector space ($n \in [1,3]$). 
The relevance of $x$ to a given
TRL level $k$ is computed via cosine similarity. The predicted label is the
index $k$ maximizing this score, provided that the maximum similarity exceeds a
noise-filtering threshold $\tau$ (set empirically to $0.15$ to avoid discarding valid but weak matches).


\subsection{Cascade Annotation Strategy for Silver Dataset Generation}

The final labeling function $\mathcal{F}(x)$ follows a priority cascade: it assigns the symbolic label $\lambda_{\text{sym}}(x)$ whenever available; otherwise, it falls back to the semantic label $\lambda_{\text{sem}}(x)$, and discards the instance if neither heuristic applies. 
This hierarchical logic guarantees that high-precision symbolic labels are never
overwritten by probabilistic semantic approximations, while preserving recall. %while still allowing the system to cover the vast majority of the corpus. 
We apply this pipeline to an extract of the CORDIS database. We use a concatenation of \emph{Title} and \emph{Objective} fields, and discard entries with
fewer than 50 characters (considered as uninformative).
From an initial dump of 35,105 projects, the pipeline successfully annotated
34,839 samples (99.2\% coverage). The symbolic heuristic
$\lambda_{\text{sym}}$ labeled 1,156 projects (3.3\%), while the semantic
heuristic $\lambda_{\text{sem}}$ labeled 33,683 projects (96.7\%). The semantic
inference achieved a mean confidence score of 0.256, and the resulting class
distribution covers the full TRL spectrum (1--9), with at least 1,400 samples
per class. We refer to this automatically annotated corpus as the
\emph{Silver Dataset}, which serves as the training set for the supervised models evaluated in Section~\ref{sec:experiments}.


\section{Experimental Setup}
\label{sec:setup}
\subsection{Two complementary corpora : silver and gold}
%GL: ajouter introduciton de la Figure fig:iea_dist en disant pourquoi la distrib des classes est différentes
As illustrated in Figure~\ref{fig:distrib_trl}, TRL levels are skewed toward early and late stages in the Silver corpus, while the Gold corpus is concentrated on higher TRLs because it focuses more strongly on technologies that are already closer to deployment and operational use.


%We rely on two complementary corpora. 
The Silver corpus is built from 34,839 CORDIS project descriptions, automatically annotated with TRL levels (1–9) by our hybrid weakly supervised pipeline described in Section \ref{sec:hybrid}. 
%that combines regular expressions and heuristic classification. 
These texts are predominantly prospective and administrative (e.g., “we will develop”, “the project aims to”), focusing on planned activities rather than established deployments. The Gold corpus consists of 508 technology notices from the IEA Clean Energy Technology Guide, manually labeled with TRLs by IEA domain experts. Notices describe relatively mature technologies in a factual, evaluative and largely retrospective style (e.g., “is widely deployed”, “has been demonstrated at commercial scale”). %There is almost no document overlap between Silver and Gold corpora. %GL: ce ne sont pas les mêmes documents de toutes façons ? C'est aps de domain overlap peut être ? Mais alors c'est déjà dit plus haut
This creates a challenging scenario where models must transfer from a mixed domain to an unseen specific one, and where the same TRL number may correspond to a shifted semantic prototype.

%MC: maths
\color{red}
Formally, we consider two data distributions:
\(
(x, y) \sim P_S \quad \text{(Silver domain)} \) and \( 
(x, y) \sim P_G \quad \text{(Gold domain)}.
\)

Models are trained on $\tilde{\mathcal{D}}_S$ sampled from $P_S$, and evaluated on $\mathcal{D}_G$ sampled from $P_G$. Domain shift is characterized by: \(P_S(x,y) \neq P_G(x,y)\).
\color{black}
%MC : maths end

%GL: pour pouvoir comparer il faudrait mettre en pourcenatge les barres car là les échelles sont trop différentes

%\begin{figure*}[t!]
%    \centering
%    \begin{subfigure}[t]{0.45\textwidth}
%  \includegraphics[width=0.95\textwidth]{silver_trl_distrib.png}
%  \caption{Silver corpus (CORDIS)}
%  \label{fig:silver_dist}
%\end{subfigure}
%~
%\begin{subfigure}[t]{0.45\textwidth}
%  \centering
%  \includegraphics[width=1\textwidth]{distriution_trl_IEA.png}
%  \caption{Gold corpus (IEA)}
%  \label{fig:iea_dist}
%\end{subfigure}

\begin{figure}[t]
    \centering
    \includegraphics[width=0.85\linewidth]{distribution_trl_concat.png}
    \caption{Distribution of TRLs in the two datasets}
    \label{fig:distrib_trl}
\end{figure}


To visualize the semantic distance between domains, we embed both corpora into a shared TF--IDF vector space and project 500 randomly sampled documents per corpus using t-SNE (after SVD to 50 dimensions). Figure~\ref{fig:tsne_domain_shift} shows two partially overlapping but clearly shifted clusters, which is consistent with a domain classifier reaching 97.8\% balanced accuracy (i.e., only $\approx$4\% effective overlap between CORDIS and IEA in the underlying semantic space).

\begin{figure}[t]%[htbp]
  \centering
  \includegraphics[width=0.5\textwidth]{fig/figure4_tsne_domain_shift.png}
  \caption{t-SNE visualization of Silver vs. Gold documents } %in a common semantic space.}
  \label{fig:tsne_domain_shift}
\end{figure}

\subsection{Models and Experimental Protocols}

We benchmark three main families of approaches, from symbolic rules to deep neural models, all aiming to recover TRLs as latent ordinal patterns from text.

\begin{enumerate}
    \item \textbf{Symbolic baseline (Regex TRL detector).} 
    We use the regex module from our labeling pipeline and apply it to the Gold corpus. This baseline quantifies how much TRL information is made explicit in the text
    and serves as a lower bound for purely symbolic reasoning.
    \item \textbf{Classical ``frugal'' models.} 
    We train classification  models (Linear SVM and Log. Regression) on TF--IDF representations  with class-balanced weights and hyperparameters ($C$, $n$-gram and vocabulary size) explored via a grid search.
    Then we train \emph{regression} models (Ridge Regression and LinearSVR) on the same representation to predict a 
    %GL: je pense que c'est trop de détails
    %continuous 
    TRL value 
    %later rounded and 
    clipped to [1, 9]. 
    Finally, we enrich this representation with grammatical cues (future tense, modals, passives), coarse TRL phases (research, proof-of-concept, demonstration, deployment), and stylometric features (sentence length, counts of future/modal/passive forms).
.    
    \item \textbf{Deep learning.} 
    We fine-tune Transformer-based models as 9-class TRL classifiers. We start with BERT, trained on the Silver corpus with a linear classification head over 9 TRL IDs. As a domain-specialised variant, we also experiment SciBERT, using the same architecture and optimisation settings. Both models are evaluated in a zero-shot Silver~$\to$~IEA transfer setting. SciBERT is further adapted in a few-shot regime on 50 labeled IEA examples. BERT consistently achieves the best zero-shot performance on the IEA corpus, while SciBERT remains slightly below despite its specialised vocabulary.
    
        
\end{enumerate}
All experiments share fixed seeds for splits and model initialization. For all symbolic and classical baselines, we
represent text with standard TF--IDF bag-of-$n$-grams, using word $n$-grams
with $n \in [1,3]$ and a maximum vocabulary size of 50{,}000 features. This
captures both single keywords and short expressions while keeping the model
size manageable. 
To address class imbalance for linear classifiers, we applied balanced class weighting during model training.
%GL: plus bas c'est spécifique à sklearn il faut être plus générique
%set \texttt{class\_weight='balanced'}
%to compensate for the skewed TRL distribution in the Silver corpus.
%\paragraph{Handling class imbalance.}
%The Silver corpus is imbalanced, with over-representation of early research (TRL~1--2) and higher deployment levels (TRL~7--9), reflecting EU funding patterns. 
%For classical models, we set \texttt{class\_weight='balanced'} to compensate for this skew. 
For deep models, we additionally apply random undersampling of majority classes during training, enforcing a maximum of $N=2,000$ examples per TRL. 

\subsubsection{Deep Learning Hyperparameters (BERT and SciBERT)}
We use the \texttt{bert-base-uncased} and %\texttt{allenai/scibert\_scivocab\_uncased} 
\texttt{scibert\_scivocab\_uncased}
checkpoints with a 9-way classification head over TRL IDs (0--8). Unless otherwise stated, both models share the same optimisation settings (learning rate $2 \cdot 10^{-5}$, weight decay $0.01$). We train for 3 epochs on the Silver corpus (batch size 16) and, for few-shot adaptation, perform a second fine-tuning stage of 10 epochs on the 50 IEA examples (batch size 8).

\subsection{Evaluation Protocol}
Due to the ordinal nature of TRLs, pure exact-match accuracy is overly strict. We report both strict accuracy and relaxed accuracy (Acc $\pm 1$), where a prediction is considered correct if $|\hat{y}-y|\le 1$. This reflects the fact that even human experts often disagree by one level when distinguishing adjacent TRLs (e.g., 4 vs 5).
We also compute {Mean Absolute Error (MAE):} The average absolute deviation $|y_i - \hat{y}_i|$, interpreted as a semantic distance on the TRL scale. We report MAE both globally and per true TRL (error profiles).
For classical models, we also compute confusion matrices (normalized by true label) on both domains, highlighting whether errors are concentrated near the diagonal (ordinal pattern) or spread across distant TRLs. For SciBERT, we additionally inspect error distributions and label-wise MAE before and after few-shot adaptation. %\paragraph{Hardware.} 
All deep learning experiments (BERT, SciBERT) are run on a single NVIDIA T4 GPU. Classical models (SVM, Logistic Regression, Naive Bayes, Random Forest, regression baselines) are trained on CPU only, illustrating the trade-off between accuracy and computational footprint.


\subsection{Transfer Protocols}
%GL: inutile je pense cette phrase ou sinon faire une vraie intro de cette subsection
%We evaluate all models under complementary training--testing regimes.

\subsubsection{Zero-Shot Transfer to the IEA Domain}

When evaluated on the Gold dataset, both frugal and deep models suffer a substantial performance drop compared to their Silver results. For linear SVMs and Logistic Regression, strict accuracy on IEA remains around 0.12--0.16 and relaxed accuracy ($\pm 1$ TRL) peaks at about 0.50--0.55, whereas on Silver the same models reach roughly 0.54 strict and 0.64--0.67 relaxed accuracy. Among deep models, we therefore use BERT base as our main reference: in a zero-shot Silver~$\to$~IEA setting it attains a strict accuracy of about 0.16 and a relaxed accuracy of about 0.58 (MAE $\approx 1.9$), clearly improving over individual frugal baselines. SciBERT, despite its domain-specialised vocabulary, performs slightly worse in this regime and does not surpass BERT base. These results support H1 on the source corpus (TRL as an ordinal latent variable), but simultaneously falsify its domain-agnostic version: even the best zero-shot models exhibit a marked degradation on IEA, with many predictions more than one level away from the expert TRL.

\paragraph{Few-Shot Domain Adaptation.}
To approximate a realistic low-resource scenario, we randomly sample $N=50$ labeled instances from the IEA corpus (stratified by TRL when possible).
For classical linear models such as SVM and Logistic Regression, we adopt a \emph{sample-boosting} strategy: the 50 IEA samples are duplicated $k$ times (typically $k=50$–$100$) and merged with the Silver corpus, effectively increasing their weight in the optimization and shifting the decision boundary toward the target domain. A similar protocol is applied to regression models (Ridge, LinearSVR), where boosted IEA instances help recalibrate continuous TRL predictions. For BERT and SciBERT, adaptation proceeds through a second fine-tuning stage in which a Silver-trained model is further trained for 10 epochs on the 50 IEA examples. On the remaining IEA documents, this few-shot stage increases strict accuracy for both models (0.33 for BERT and 0.31 for SciBERT) and improves macro-F1, while preserving relaxed accuracy 0.56 and 0.52, respectively. Overall, BERT few-shot provides the strongest deep-learning approximation of the TRL latent pattern on the target domain, but both deep models still remain equivalent to the best frugal hybrids in relaxed accuracy.

Together, these two protocols disentangle (i) how well different architectures learn the TRL scale as an ordinal pattern in-domain, and (ii) how sensitive each family of models is to semantic domain shift between administrative project descriptions and expert technology assessments.

\section{Research hypothesis and experiments}
\label{sec:hyp}
\label{sec:experiments}

The overarching goal of this work is to identify what conveys information about technological maturity in text. Expert practice suggests that TRL levels are rarely stated explicitly, but emerge as a \emph{latent pattern} distributed across vocabulary, sentence structure and discourse mode. We posit that this signal can be modeled (i) as an \emph{ordinal latent variable} and (ii) as a \emph{partially grammaticalized structure}. These two properties guide our experimental study.

\newtheorem{hypothesis}{Hypothesis}
\begin{hypothesis}[H1. Ordinal latent variable]\label{H1}
The TRL scale behaves as a continuous and ordered latent variable encoded in textual descriptions.
If this latent maturity pattern is truly domain-agnostic, models trained on one scientific domain should recover the same ordinal structure (i.e., predominantly $\pm 1$ errors) when transferred to another domain.\end{hypothesis} 

\begin{hypothesis}[H2. Grammatical structure]\label{H2}
The latent TRL pattern is associated with stable grammatical cues such as tense, modality, voice, and phase markers which formalize the progression of technological maturity.
If these cues drive TRL modelling, representations that explicitly encode grammatical structure should improve prediction quality and robustness.
\end{hypothesis}

In order to test these two hypotheses we first examine whether TRLs behave as an
ordinal latent variable that can be learned on the Silver corpus and transferred
to the IEA domain (H1). We then investigate the role of grammatical cues as a
possible formalization of technological maturity (H2).

\subsection{Experiment 1: Ordinal TRL Pattern and Domain Shift (H1)}
\label{sec:exp1}

\subsubsection{Frugal Linear Models on the Silver Corpus}
Using TF--IDF representations, we benchmark linear SVMs and Logistic Regression with class-balanced weights
and varied regularization $C$. 
On the Silver corpus, the best configuration is a linear SVM (with $C = 0.1$ and $n$-grams 1--2) which achieves around 55\% strict accuracy and 66\%
relaxed accuracy .
This confirms that TRL levels indeed behave as
an ordered latent variable on the source domain: most errors remain within one
level of the true label.
To better understand how the model organizes the TRL scale, we inspect the top
SVM coefficients for selected levels. Figure~\ref{fig:topwords_silver} shows
that TRL~1 is dominated by research-related terms, TRL~4 by laboratory and
controlled-environment vocabulary, TRL~6 by pilot and relevant-environment
phrases, and TRL~9 by market and deployment markers, providing qualitative
evidence for an ordered latent pattern and remains consistent with the original TRL scale.

\begin{figure}[t]%[h!]
  \centering
  \includegraphics[width=0.7\textwidth]{ICPR_2026_LaTeX_Templates/figure_feature_importance_silver.png}
  \caption{Top lexical predictors for selected TRL levels (1, 4, 6, 9) learned by
  the SVM on the Silver corpus.}
  \label{fig:topwords_silver}
\end{figure}
%GL: TODO traduire les termes et si possible standardiser l'échelle en abscisse

\subsubsection{Zero-Shot Transfer to the IEA Domain}

When evaluated on the Gold dataset, both frugal and deep models suffer a substantial performance drop compared to their Silver results. For linear SVMs and Logistic Regression, strict accuracy on IEA remains around 0.12--0.16 and relaxed accuracy ($\pm 1$ TRL) peaks at about 0.50--0.55, whereas on Silver the same models reach roughly 0.54 strict and 0.64--0.67 relaxed accuracy. Among deep models, we therefore use BERT base as our reference: in a zero-shot Silver~$\to$~IEA setting it attains a strict accuracy of 0.14 and a relaxed accuracy of 0.58 (MAE: 1.92), clearly improving over the frugal baselines. On the Silver corpus, SciBERT slightly outperforms BERT (strict accuracy: 0.66, relaxed accuracy: 0.74, MAE: 1.21), but under zero-shot transfer to IEA it reaches a strict accuracy of 0.15 and a relaxed accuracy of 0.54 with a higher MAE (2.03), and is thus not retained as our main deep reference. These results support H1 on the source corpus (TRL as an ordinal latent variable), but simultaneously falsify its domain-agnostic version: even the best zero-shot models exhibit a marked degradation on IEA, with many predictions more than one level away from the expert TRL.

\subsubsection{Ordinal Regression vs. Multi-class Classification}

To explicitly test the ordinal nature of TRLs, we compare regression and
classification models trained on Silver and evaluated on IEA.
Regression models achieve slightly lower MAE (e.g., 2.28 for Linear SVR vs.\ 2.47
for SVM), confirming that their errors tend to be closer to the true TRL. However, they underperform in relaxed accuracy: Linear SVR trails the
SVM by almost 8 percentage points. Error histograms reveal that regression
outputs form a smooth distribution around the true labels but lack the sharpness
required for precise decision boundaries between adjacent TRLs. Certain lexical
or discourse cues (e.g., explicit ``pilot plant'', ``operational environment'')
act as quasi-categorical triggers, better captured by multi-class
classification.
Overall, these findings confirm that TRLs are organized as an ordinal latent
variable (errors are mostly grouped around matrix diagonal), but
multi-class classification remains preferable for practical TRL assignment, as
it learns distinct semantic prototypes per level while preserving an ordinal
error structure.


\subsubsection{Few-Shot Realignment and Specialized Language Models}

To test whether a small amount of IEA supervision can realign the TRL pattern, we simulate few-shot adaptation with 50 %$N=50$ 
labeled IEA examples. Few-shot transfer on classical models yields consistent improvements on IEA relative to pure Silver-only training (higher relaxed accuracy, lower MAE), but the absolute performance remains limited. For BERT base, we keep the zero-shot setting as our main deep reference (strict accuracy $\approx 0.14$, relaxed accuracy $\approx 0.58$, MAE $\approx 1.92$). SciBERT provides a complementary perspective: in zero-shot transfer it remains slightly below BERT in relaxed accuracy on IEA (0.54 vs. 0.58) and exhibits a higher MAE (2.03), but few-shot fine-tuning raises its strict accuracy to about 0.31 and improves macro-F1, with relaxed accuracy of 0.52 and MAE: 1.84. Overall, these experiments confirm that deep contextual models can learn a coherent ordinal TRL pattern and benefit from limited target supervision, but also that domain-specialised vocabularies (SciBERT) do not automatically outperform a well-tuned general model (BERT base) in this setting.
Table~\ref{tab:deep_trl_iea} summarizes the performance of the deep models on the IEA corpus in zero-shot and few-shot settings.



\begin{table}[b]
\centering
\small
\begin{tabular}{l l l c c c c}
\toprule
Model & Scenario & Test set & Strict Acc. & Acc. ($\pm 1$) & F1$_\text{macro}$ & MAE \\
\midrule
BERT (base) & Zero-shot & Silver & 61.6\% & 69.9\% & 0.422 & 1.37 \\
BERT (base) & Zero-shot & IEA    & 13.7\% & 57.6\% & 0.100 & 1.92 \\
BERT (base) & Few-shot  & Silver & 44.2\% & 65.7\% & 0.347 & 1.69 \\
BERT (base) & Few-shot  & IEA    & 33.4\% & 56.0\% & 0.172 & 1.70 \\
\midrule
SciBERT      & Zero-shot & Silver & 66.0\% & 73.6\% & 0.511 & 1.21 \\
SciBERT      & Zero-shot & IEA    & 14.8\% & 54.3\% & 0.096 & 2.03 \\
SciBERT      & Few-shot  & Silver & 46.4\% & 67.6\% & 0.406 & 1.61 \\
SciBERT      & Few-shot  & IEA    & 31.0\% & 51.8\% & 0.166 & 1.84 \\
\bottomrule
\end{tabular}
\caption{Deep learning models for TRL prediction on Silver and IEA corpora, under zero-shot and few-shot settings.}
\label{tab:deep_trl_iea}
\end{table}

\subsection{Experiment 2: Grammatical Cues as a Formalization of Maturity (H2)}
\label{sec:exp2}

\subsubsection{Grammar-based Representations}
\label{gramm}
We test whether explicit grammatical structure helps model and transfer the TRL
pattern (H2). We design four alternative representations, evaluated under a
unified protocol: training on 80\% of Silver, internal testing on the remaining
20\%, and external evaluation on Gold.

\begin{description}
    \item[(R1) Full-text TF--IDF:] word $n$-grams (1--3) for the raw text.
    \item[(R2) Grammar+TRL view:] stylistic tags (future, modal, passive) and
    coarse phase markers (e.g., \emph{research}, \emph{proof of concept},
    \emph{demonstration}, \emph{deployment}), concatenated with TRL terms derived
    from extended definitions.
    \item[(R3) Structural features:] numeric stylometric indicators (tense,
    modality, passives, phase markers, text length).
    \item[(R4) Fusion:] concatenation of R1 and R3.
\end{description}

For each representation, we train linear SVMs, Logistic Regression, and Random
Forests, and report strict accuracy, relaxed accuracy, MAE, and the proportion
of catastrophic errors ($|\hat{y}-y|\ge3$).

\subsubsection{Results on the Silver Corpus}

On the Silver corpus, the Grammar+TRL view (R2) yields modest but consistent improvements over full-text (R1) for linear models. For example, with a linear SVM, relaxed accuracy increases from 0.65 to 0.69, while MAE decreases from 1.53 to 1.32 and catastrophic errors from 0.26 to 0.23. This indicates that grammatical
scaffolding combined with TRL lexemes captures part of the maturity structure present in weak labels. In contrast, purely structural features (R3) perform poorly (strict accuracy
$\approx 0.13$--$0.17$, MAE $\approx 3.0$--$3.3$), showing that grammar alone,
without lexical content, is insufficient to recover fine-grained ordinal
progression. Fusion (R4) does not yield systematic gains beyond R1 or R2.

\subsubsection{Generalization to the Gold Domain}

Under transfer to the IEA Gold corpus, the advantages of R2 disappear or reverse.
For the same SVM, relaxed accuracy drops from 0.51 (R1) to 0.44 (R2), while MAE
and catastrophic errors increase (2.10 to 2.54 and 0.37 to 0.43, respectively).
Similar trends are observed for Logistic Regression and Random Forest. Neither
Grammar+TRL nor fusion representations consistently improve ordinal robustness
on Gold, indicating that grammar-centric views do not mitigate the Silver-to-IEA
domain shift.

\subsubsection{Error Biases and Linguistic Motifs}

To further analyze this mismatch, we examine the most influential $n$-grams
driving SVM predictions on IEA. Figure~\ref{fig:triggerwords_gold} shows that
high-TRL triggers are highly domain-specific (e.g., \emph{algae},
\emph{methanol fuel}, \emph{PEM fuel}, \emph{clinker}), and differ markedly from
Silver prototypes. This highlights that the semantic realization of TRLs is not
stable across domains.



\begin{figure}[t]
  \centering
  \begin{subfigure}[t]{0.55\textwidth}
    \centering
    \includegraphics[width=\textwidth]{ICPR_2026_LaTeX_Templates/figure_prediction_drivers_gold.png}
    \caption{Trigger words for TRL~6--9 (IEA).}
  \end{subfigure}\hfill
  \begin{subfigure}[t]{0.40\textwidth}
    \centering
    \includegraphics[width=\textwidth]{fig_bias_patterns_full_vs_grammar.png}
    \caption{Average prediction errors $(\hat{y}-y)$ for key motifs.}
  \end{subfigure}
  \caption{Error drivers and bias patterns for TRL prediction on the IEA corpus.}
  \label{fig:gold_drivers_bias}
\end{figure}

We further compare prediction biases for key linguistic motifs under full-text
(R1) and Grammar+TRL (R2) representations. For each motif (future, modal,
research, demonstration, deployment), we compute the mean error $(\hat{y}-y)$ on
IEA. Both models exhibit nearly identical bias profiles: they systematically
underestimate TRLs for future-, modal-, research-, and deployment-related
expressions (about $-0.5$ to $-1.2$ levels), and slightly overestimate
demonstration-related terms. Grammar+TRL modifies error magnitudes marginally
but does not change bias direction or interpretation
(Fig.~\ref{fig:gold_drivers_bias}). Taken together, these results falsify the strong form of Hypothesis~\ref{H2}.
While grammatical cues help interpret what models learn in-domain, they do not
provide a domain-agnostic formalization of maturity. The dominant source of
domain shift between Silver and IEA is therefore semantic and operational rather
than grammatical, reflecting differences in how maturity-related terms are
mapped to TRL levels across institutions.

\section{Discussion and Conclusion}
%strict vs fuzzy
% TABLEAU ALL RECAP
\input{ICPR_2026_LaTeX_Templates/recap_tab}
Table~\ref{tab:recap_tab} summarizes the best-performing models per methodological family under the minimized Silver supervision scheme. In-domain results show that weak labels support a structured ordinal interpretation, with hard hierarchical models achieving the lowest ordinal errors. Under domain shift to the IEA corpus, however, rigid ordinal constraints degrade, while soft hierarchical routing significantly improves robustness. The best frugal performance is obtained by a hybrid model combining an ordinal expert and a local pattern aggregator, indicating that TRL maturity reflects both a global ordered structure and context-specific decision rules. Strict accuracy remains uniformly low under domain shift for all frugal models, confirming that it is not a reliable indicator of ordinal robustness; we therefore emphasize relaxed accuracy ($\pm 1$), MAE, and catastrophic errors.

Overall, our experiments confirm that TRLs behave as an ordered latent pattern in text, but that this pattern is strongly shaped by domain conventions. On the Silver corpus derived from CORDIS, frugal linear models trained on weakly supervised labels recover a coherent ordinal structure: errors concentrate within $\pm 1$ level and class analysis reveals intuitive prototypes for early (research), intermediate (laboratory and pilot), and late (deployment and market) stages. This supports Hypothesis~\ref{H1} in the source domain and shows that weak supervision can enable accurate and computationally efficient TRL modelling at scale.

However, the same models degrade under zero-shot transfer to the IEA Gold corpus, despite relying on identical textual inputs and a shared TRL scale. Confusion matrices and error profiles reveal a substantial semantic drift between CORDIS project summaries and IEA technology notes. Terms such as ``pilot plant'', ``demonstration'' or ``deployment'' systematically correspond to higher TRL levels in IEA than in CORDIS, even when their surface form is similar. This contradicts the strong, domain-agnostic version of Hypothesis~\ref{H1}: ordinal maturity patterns are not stable across domains, and models trained on administrative data cannot be deployed unchanged to expert corpora.

Grammar-aware representations provide additional insight but do not resolve this mismatch. While grammatical scaffolding and coarse phase markers improve performance on the Silver corpus and help interpret progression cues, they do not improve relaxed accuracy, MAE, or catastrophic errors on IEA. Bias analyses by linguistic motif show that full-text and Grammar+TRL models exhibit nearly identical error profiles, indicating that domain shift is driven primarily by semantic and operational conventions rather than by syntactic form. This falsifies the strong form of Hypothesis~\ref{H2}: grammatical cues alone do not provide a domain-agnostic formalisation of TRLs.

From an operational perspective, these results suggest a hybrid strategy for large-scale TRL assessment. Weak supervision offers a scalable way to construct high-coverage Silver datasets from administrative sources, enabling efficient in-domain scoring with frugal linear models. For cross-domain deployment, however, additional flexibility or supervision is required. Our results show that hybrid frugal models combining ordinal structure and local pattern aggregation substantially reduce large ordinal errors under domain shift, and that a small amount of expert annotation can further realign maturity estimates when coupled with richer language models such as BERT (and, more generally, Transformer-based encoders). 
\color{red}


%MC : limitations
\paragraph{Limitations.}
This study presents several limitations. First, the Silver dataset relies on weak supervision and may contain systematic biases induced by the labeling heuristics. Second, the domain shift analysis is limited to a single target domain (energy technologies), which may not generalize to other sectors. Third, the evaluation focuses on textual signals and does not take in consideration structure of documents neither metadata, which may carry complementary information about technological maturity.
%MC : fin limitations
\color{black}
\paragraph{Future work.}
Several avenues follow from this study. First, further work is needed to characterise how experts in different organisations map textual evidence to TRL levels. Comparative annotation studies and protocol analyses could help formalise implicit decision rules and improve the design of weak labeling functions. Second, extending the framework to additional domains and innovation ecosystems would allow systematic measurement of how TRL semantics vary across contexts. Third, richer representations of maturity—combining refined sector-specific lexicons with structured project information (e.g., milestones, work packages, performance indicators) offer a promising direction for more robust and interpretable TRL estimation across heterogeneous data sources.

Overall, this work shows that while TRLs manifest as an ordered latent pattern that can be learned from weak supervision, this pattern is not directly portable across domains. Scalable and robust TRL assessment therefore requires a combination of weakly supervised breadth, ordinal inductive bias, and domain-aware calibration, rather than a single domain-agnostic classifier.
All datasets and scripts used in this study are publicly available on GitHub to support reproducibility and reuse\footnote{\url{https://github.com/melklr/ICPR\_2026}}.
%GL: la footnote était après le  point final de la phrase, j'ai changé.




%\newpage
%
% ---- Bibliography ----
%
% BibTeX users should specify bibliography style 'splncs04'.
% References will then be sorted and formatted in the correct style.
%
\bibliographystyle{splncs04}
\bibliography{ICPR_2026_LaTeX_Templates/biblio}
%

\end{document}
